% ============================================================
% KI-Nutzungserklärung (nach OWL-Variante 1)
% Einzufügen NACH dem Literaturverzeichnis / im Anhang.
% Einbindung z. B. via % ============================================================
% KI-Nutzungserklärung (nach OWL-Variante 1)
% Einzufügen NACH dem Literaturverzeichnis / im Anhang.
% Einbindung z. B. via % ============================================================
% KI-Nutzungserklärung (nach OWL-Variante 1)
% Einzufügen NACH dem Literaturverzeichnis / im Anhang.
% Einbindung z. B. via % ============================================================
% KI-Nutzungserklärung (nach OWL-Variante 1)
% Einzufügen NACH dem Literaturverzeichnis / im Anhang.
% Einbindung z. B. via \include{content/ki_erklaerung} im Anhang-Bereich.
% Grundlage: KI-Leitlinie der TH Nürnberg (Amtsblatt 2/2026) und
% OWL-Handreichung "KI-Nutzung in wissenschaftlichen Arbeiten dokumentieren".
% ============================================================

\chapter*{Erklärung: Verwendung von KI-Werkzeugen}
\addcontentsline{toc}{chapter}{Erklärung: Verwendung von KI-Werkzeugen}

Im Rahmen der Erstellung dieser Arbeit wurden generative KI-Werkzeuge eingesetzt.
Die Nutzung wird im Folgenden gemäß der KI-Leitlinie der Technischen Hochschule
Nürnberg Georg Simon Ohm offengelegt.

\begin{itemize}
  \item \textbf{Claude Opus 4.8 (Anthropic):}
        Unterstützung bei der Strukturierung und Gliederung der Arbeit,
        bei der Literaturrecherche mit anschließender Verifikation aller
        Quellen an den Originalpublikationen,
        bei der Erstellung und Überarbeitung von Textentwürfen
        sowie als dialogisches Gegenüber zur Schärfung der Argumentation
        und zur Reflexion des methodischen Vorgehens.
\end{itemize}

Sämtliche von KI-Werkzeugen erzeugten Inhalte wurden eigenständig geprüft,
überarbeitet und inhaltlich verantwortet.
Die Verantwortung für alle Aussagen, Ergebnisse und Schlussfolgerungen dieser
Arbeit liegt vollständig beim Verfasser.
Die reine Rechtschreib- und Grammatikprüfung ist gemäß Leitlinie nicht
dokumentationspflichtig und daher nicht gesondert aufgeführt.

\bigskip
\noindent\textbf{Dokumentationspflichtiger Prompt (inhaltlich prägend).}
Der folgende Prompt trug maßgeblich zur Fundierung der Forschungslücke
(Kapitel~\ref{ch:einleitung} und~\ref{ch:grundlagen}) bei:

\begin{quote}
\enquote{Führe eine wissenschaftliche Literaturrecherche zu Szenenrepräsentationen
für LLM/VLM-Agenten bei natürlichsprachlicher 3D-Szenenmanipulation durch.
Sammle Belege für drei Stränge --- strukturelle, visuelle und hybride
Repräsentation --- sowie für die dazwischen bestehende Forschungslücke,
und schätze ab, wie gut die Lückenbehauptung
(\enquote{ein systematischer, referenztyp-differenzierter Vergleich struktureller,
visueller und hybrider Repräsentation für lokale Agenten fehlt}) gestützt wird.}
\end{quote}

\noindent
Claude Opus 4.8, Anthropic. Die durch die Recherche vorgeschlagenen Quellen
wurden ausnahmslos an den Originalpublikationen überprüft; nicht verifizierbare
Angaben wurden verworfen.
 im Anhang-Bereich.
% Grundlage: KI-Leitlinie der TH Nürnberg (Amtsblatt 2/2026) und
% OWL-Handreichung "KI-Nutzung in wissenschaftlichen Arbeiten dokumentieren".
% ============================================================

\chapter*{Erklärung: Verwendung von KI-Werkzeugen}
\addcontentsline{toc}{chapter}{Erklärung: Verwendung von KI-Werkzeugen}

Im Rahmen der Erstellung dieser Arbeit wurden generative KI-Werkzeuge eingesetzt.
Die Nutzung wird im Folgenden gemäß der KI-Leitlinie der Technischen Hochschule
Nürnberg Georg Simon Ohm offengelegt.

\begin{itemize}
  \item \textbf{Claude Opus 4.8 (Anthropic):}
        Unterstützung bei der Strukturierung und Gliederung der Arbeit,
        bei der Literaturrecherche mit anschließender Verifikation aller
        Quellen an den Originalpublikationen,
        bei der Erstellung und Überarbeitung von Textentwürfen
        sowie als dialogisches Gegenüber zur Schärfung der Argumentation
        und zur Reflexion des methodischen Vorgehens.
\end{itemize}

Sämtliche von KI-Werkzeugen erzeugten Inhalte wurden eigenständig geprüft,
überarbeitet und inhaltlich verantwortet.
Die Verantwortung für alle Aussagen, Ergebnisse und Schlussfolgerungen dieser
Arbeit liegt vollständig beim Verfasser.
Die reine Rechtschreib- und Grammatikprüfung ist gemäß Leitlinie nicht
dokumentationspflichtig und daher nicht gesondert aufgeführt.

\bigskip
\noindent\textbf{Dokumentationspflichtiger Prompt (inhaltlich prägend).}
Der folgende Prompt trug maßgeblich zur Fundierung der Forschungslücke
(Kapitel~\ref{ch:einleitung} und~\ref{ch:grundlagen}) bei:

\begin{quote}
\enquote{Führe eine wissenschaftliche Literaturrecherche zu Szenenrepräsentationen
für LLM/VLM-Agenten bei natürlichsprachlicher 3D-Szenenmanipulation durch.
Sammle Belege für drei Stränge --- strukturelle, visuelle und hybride
Repräsentation --- sowie für die dazwischen bestehende Forschungslücke,
und schätze ab, wie gut die Lückenbehauptung
(\enquote{ein systematischer, referenztyp-differenzierter Vergleich struktureller,
visueller und hybrider Repräsentation für lokale Agenten fehlt}) gestützt wird.}
\end{quote}

\noindent
Claude Opus 4.8, Anthropic. Die durch die Recherche vorgeschlagenen Quellen
wurden ausnahmslos an den Originalpublikationen überprüft; nicht verifizierbare
Angaben wurden verworfen.
 im Anhang-Bereich.
% Grundlage: KI-Leitlinie der TH Nürnberg (Amtsblatt 2/2026) und
% OWL-Handreichung "KI-Nutzung in wissenschaftlichen Arbeiten dokumentieren".
% ============================================================

\chapter*{Erklärung: Verwendung von KI-Werkzeugen}
\addcontentsline{toc}{chapter}{Erklärung: Verwendung von KI-Werkzeugen}

Im Rahmen der Erstellung dieser Arbeit wurden generative KI-Werkzeuge eingesetzt.
Die Nutzung wird im Folgenden gemäß der KI-Leitlinie der Technischen Hochschule
Nürnberg Georg Simon Ohm offengelegt.

\begin{itemize}
  \item \textbf{Claude Opus 4.8 (Anthropic):}
        Unterstützung bei der Strukturierung und Gliederung der Arbeit,
        bei der Literaturrecherche mit anschließender Verifikation aller
        Quellen an den Originalpublikationen,
        bei der Erstellung und Überarbeitung von Textentwürfen
        sowie als dialogisches Gegenüber zur Schärfung der Argumentation
        und zur Reflexion des methodischen Vorgehens.
\end{itemize}

Sämtliche von KI-Werkzeugen erzeugten Inhalte wurden eigenständig geprüft,
überarbeitet und inhaltlich verantwortet.
Die Verantwortung für alle Aussagen, Ergebnisse und Schlussfolgerungen dieser
Arbeit liegt vollständig beim Verfasser.
Die reine Rechtschreib- und Grammatikprüfung ist gemäß Leitlinie nicht
dokumentationspflichtig und daher nicht gesondert aufgeführt.

\bigskip
\noindent\textbf{Dokumentationspflichtiger Prompt (inhaltlich prägend).}
Der folgende Prompt trug maßgeblich zur Fundierung der Forschungslücke
(Kapitel~\ref{ch:einleitung} und~\ref{ch:grundlagen}) bei:

\begin{quote}
\enquote{Führe eine wissenschaftliche Literaturrecherche zu Szenenrepräsentationen
für LLM/VLM-Agenten bei natürlichsprachlicher 3D-Szenenmanipulation durch.
Sammle Belege für drei Stränge --- strukturelle, visuelle und hybride
Repräsentation --- sowie für die dazwischen bestehende Forschungslücke,
und schätze ab, wie gut die Lückenbehauptung
(\enquote{ein systematischer, referenztyp-differenzierter Vergleich struktureller,
visueller und hybrider Repräsentation für lokale Agenten fehlt}) gestützt wird.}
\end{quote}

\noindent
Claude Opus 4.8, Anthropic. Die durch die Recherche vorgeschlagenen Quellen
wurden ausnahmslos an den Originalpublikationen überprüft; nicht verifizierbare
Angaben wurden verworfen.
 im Anhang-Bereich.
% Grundlage: KI-Leitlinie der TH Nürnberg (Amtsblatt 2/2026) und
% OWL-Handreichung "KI-Nutzung in wissenschaftlichen Arbeiten dokumentieren".
% ============================================================

\chapter*{Erklärung: Verwendung von KI-Werkzeugen}
\addcontentsline{toc}{chapter}{Erklärung: Verwendung von KI-Werkzeugen}

Im Rahmen der Erstellung dieser Arbeit wurden generative KI-Werkzeuge eingesetzt.
Die Nutzung wird im Folgenden gemäß der KI-Leitlinie der Technischen Hochschule
Nürnberg Georg Simon Ohm offengelegt.

\begin{itemize}
  \item \textbf{Claude Opus 4.8 (Anthropic):}
        Unterstützung bei der Strukturierung und Gliederung der Arbeit,
        bei der Literaturrecherche mit anschließender Verifikation aller
        Quellen an den Originalpublikationen,
        bei der Erstellung und Überarbeitung von Textentwürfen
        sowie als dialogisches Gegenüber zur Schärfung der Argumentation
        und zur Reflexion des methodischen Vorgehens.
\end{itemize}

Sämtliche von KI-Werkzeugen erzeugten Inhalte wurden eigenständig geprüft,
überarbeitet und inhaltlich verantwortet.
Die Verantwortung für alle Aussagen, Ergebnisse und Schlussfolgerungen dieser
Arbeit liegt vollständig beim Verfasser.
Die reine Rechtschreib- und Grammatikprüfung ist gemäß Leitlinie nicht
dokumentationspflichtig und daher nicht gesondert aufgeführt.

\bigskip
\noindent\textbf{Dokumentationspflichtiger Prompt (inhaltlich prägend).}
Der folgende Prompt trug maßgeblich zur Fundierung der Forschungslücke
(Kapitel~\ref{ch:einleitung} und~\ref{ch:grundlagen}) bei:

\begin{quote}
\enquote{Führe eine wissenschaftliche Literaturrecherche zu Szenenrepräsentationen
für LLM/VLM-Agenten bei natürlichsprachlicher 3D-Szenenmanipulation durch.
Sammle Belege für drei Stränge --- strukturelle, visuelle und hybride
Repräsentation --- sowie für die dazwischen bestehende Forschungslücke,
und schätze ab, wie gut die Lückenbehauptung
(\enquote{ein systematischer, referenztyp-differenzierter Vergleich struktureller,
visueller und hybrider Repräsentation für lokale Agenten fehlt}) gestützt wird.}
\end{quote}

\noindent
Claude Opus 4.8, Anthropic. Die durch die Recherche vorgeschlagenen Quellen
wurden ausnahmslos an den Originalpublikationen überprüft; nicht verifizierbare
Angaben wurden verworfen.
